\documentclass[UTF8,zihao=-4]{ctexart}
\usepackage[a4paper,margin=2.2cm]{geometry}
\usepackage{xcolor}
\usepackage{hyperref}
\usepackage{enumitem}
\usepackage{fontspec}
\usepackage{amsmath}
\IfFontExistsTF{Microsoft YaHei}{\setCJKmainfont{Microsoft YaHei}}{\setCJKmainfont{SimSun}[BoldFont=SimHei]}
\IfFontExistsTF{Microsoft YaHei UI}{\setCJKsansfont{Microsoft YaHei UI}}{}
\IfFontExistsTF{Consolas}{\setmonofont{Consolas}}{}
\definecolor{linkblue}{RGB}{30,80,145}
\hypersetup{colorlinks=true,linkcolor=linkblue,urlcolor=linkblue,pdfborder={0 0 0}}
\setlength{\parindent}{2em}
\setlength{\parskip}{0.25em}
\linespread{1.16}
\setlist{itemsep=0.2em,topsep=0.35em,leftmargin=2.5em}
\pagestyle{plain}

\title{10\_DFS与BFS图搜索}

\author{}

\date{}

\begin{document}

\maketitle

\section*{10 DFS 与 BFS 图搜索}

\subsection*{题目原文}

简述深度优先和广度优先算法的思想，采用该算法思想能解决经典的图上搜索问题：topo 序、任务调度、关键路径和强联通区域问题

\subsection*{考查类型}

概念题和算法设计题。

\subsection*{DFS 思想}

深度优先搜索从一个起点出发，沿着一条路径尽量向深处走；当没有未访问邻接点时，再回退到上一个结点继续搜索。

实现方式：

\begin{quote}
\small\ttfamily\raggedright
\relax 递归\\
\relax 显式栈\\
\end{quote}

典型记录：

\begin{quote}
\small\ttfamily\raggedright
\relax visited[v]\\
\relax discover\ time\\
\relax finish\ time\\
\relax parent[v]\\
\end{quote}

时间复杂度：

\begin{quote}
\small\ttfamily\raggedright
\relax O(V\ +\ E)\\
\end{quote}

\subsection*{BFS 思想}

广度优先搜索从起点出发，先访问距离为 1 的点，再访问距离为 2 的点，按层推进。

实现方式：

\begin{quote}
\small\ttfamily\raggedright
\relax 队列\\
\end{quote}

典型记录：

\begin{quote}
\small\ttfamily\raggedright
\relax visited[v]\\
\relax distance[v]\\
\relax parent[v]\\
\end{quote}

时间复杂度：

\begin{quote}
\small\ttfamily\raggedright
\relax O(V\ +\ E)\\
\end{quote}

BFS 可求无权图中从起点到各点的最短路径长度。

\subsection*{topo 序}

拓扑序用于有向无环图。

DFS 做法：

\begin{enumerate}
\item 对图做 DFS。
\item 按完成时间从大到小排列顶点。
\item 得到拓扑序。
\end{enumerate}

Kahn 算法做法：

\begin{enumerate}
\item 统计所有顶点入度。
\item 将入度为 0 的顶点入队。
\item 每取出一个顶点，就删除它的出边并更新邻点入度。
\item 新的入度为 0 的顶点入队。
\end{enumerate}

若最终输出顶点数少于 \texttt{V}，说明图中存在环。

\subsection*{任务调度}

任务依赖关系可建模为有向图：

\begin{quote}
\small\ttfamily\raggedright
\relax u\ ->\ v\ 表示任务\ u\ 必须先于任务\ v\ 完成\\
\end{quote}

若图是 DAG，则拓扑序给出合法任务执行顺序。

若图有环，则依赖关系冲突，无法完成全部任务。

\subsection*{关键路径}

关键路径通常在有向无环图中讨论。边或点带有持续时间。

基本步骤：

\begin{enumerate}
\item 求拓扑序。
\item 按拓扑序计算每个点的最早开始或最早完成时间。
\item 反向按逆拓扑序计算最晚开始或最晚完成时间。
\item 松弛时间为 0 的活动位于关键路径上。
\end{enumerate}

关键路径长度决定项目最短完成时间。

\subsection*{强连通区域问题}

有向图中，若两个顶点互相可达，则它们属于同一个强连通分量。

Kosaraju 算法：

\begin{enumerate}
\item 对原图做 DFS，记录完成时间。
\item 将所有边反向。
\item 按原图完成时间从大到小，在反图上 DFS。
\item 每次 DFS 访问到的一组顶点构成一个强连通分量。
\end{enumerate}

Tarjan 算法也使用 DFS，通过 \texttt{dfn} 和 \texttt{low} 维护强连通分量。

\subsection*{例题}

\subsubsection*{例题 1}

题目：给定有向图：

\[V = {A, B, C, D, E, F}\]
\[E = {A->C, B->C, B->D, C->E, D->F, E->F}\]

写出一个拓扑序。

解答：

一个合法拓扑序为：

\begin{quote}
\small\ttfamily\raggedright
\relax A,\ B,\ C,\ D,\ E,\ F\\
\end{quote}

检查每条边：

\begin{quote}
\small\ttfamily\raggedright
\relax A\ 在\ C\ 前\\
\relax B\ 在\ C\ 前\\
\relax B\ 在\ D\ 前\\
\relax C\ 在\ E\ 前\\
\relax D\ 在\ F\ 前\\
\relax E\ 在\ F\ 前\\
\end{quote}

所有依赖均满足，因此该序列是拓扑序。

\subsubsection*{例题 2}

题目：在无权图中从 \texttt{s} 出发做 BFS。边集合为：

\begin{quote}
\small\ttfamily\raggedright
\relax s-a,\ s-b,\ a-c,\ b-c,\ c-d,\ b-e\\
\end{quote}

求从 \texttt{s} 到各点的最短距离。

解答：

BFS 按层推进：

\begin{quote}
\small\ttfamily\raggedright
\relax 第\ 0\ 层：s\\
\relax 第\ 1\ 层：a,\ b\\
\relax 第\ 2\ 层：c,\ e\\
\relax 第\ 3\ 层：d\\
\end{quote}

距离为：

\begin{quote}
\small\ttfamily\raggedright
\relax dist[s]\ =\ 0\\
\relax dist[a]\ =\ 1\\
\relax dist[b]\ =\ 1\\
\relax dist[c]\ =\ 2\\
\relax dist[e]\ =\ 2\\
\relax dist[d]\ =\ 3\\
\end{quote}

\subsubsection*{例题 3}

题目：求下面有向图的强连通分量：

\[1->2, 2->3, 3->1, 3->4, 4->5, 5->4\]

解答：

顶点 \texttt{1, 2, 3} 互相可达，构成一个强连通分量：

\begin{quote}
\small\ttfamily\raggedright
\relax \{1,\ 2,\ 3\}\\
\end{quote}

顶点 \texttt{4, 5} 互相可达，构成一个强连通分量：

\begin{quote}
\small\ttfamily\raggedright
\relax \{4,\ 5\}\\
\end{quote}

边 $3->4$ 只提供从第一个分量到第二个分量的方向，不能让两个分量合并。

\subsection*{常见误区}

BFS 适合无权图最短路，不适合直接处理带权最短路。

拓扑序只对 DAG 存在。若图有有向环，则不存在拓扑序。

强连通分量是有向图概念，不是无向连通分量。

\end{document}
